\documentclass[11pt,a4paper]{article}
\usepackage[margin=2.5cm]{geometry}
\usepackage[T1]{fontenc}
\usepackage[utf8]{inputenc}
\usepackage{lmodern}
\usepackage{microtype}
\usepackage{booktabs}
\usepackage{enumitem}
\usepackage{xcolor}
\usepackage{listings}
\usepackage[hidelinks]{hyperref}
\setlist{nosep}
\lstset{basicstyle=\ttfamily\small,breaklines=true,frame=single,
  columns=fullflexible}

\title{\textbf{DBnomics for Macroeconomic Data}\\
\large A reproducible guide for teaching, MATLAB, Dynare, and Python}
\author{Gauthier Vermandel}
\date{21 August 2026}

\begin{document}
\maketitle

\noindent\textbf{Repository.} \url{https://github.com/Vermandel/dbnomics-teaching}\\
\textbf{Purpose.} This is a generic companion guide for courses using
macroeconomic data. It explains how to find, download, document, and prepare
series from DBnomics. It is not specific to one assignment or one model.

\section{The reproducible workflow}

DBnomics makes datasets from many providers available through a common web
interface and API. It does not remove the need to understand the original
source, definitions, units, seasonal adjustment, frequency, licence, or
revision policy. The reliable sequence is:

\begin{enumerate}
\item Find a series in the DBnomics web interface and inspect its metadata.
\item Record its \emph{provider}, \emph{dataset}, and \emph{series code}.
\item Call the API from a script, rather than committing a downloaded file.
\item Preserve the raw period label and missing values; transform them
      explicitly and document the choice.
\item Save a small, model-ready CSV only locally or as a build product.
\end{enumerate}

An API endpoint has the form
\begin{center}
\url{https://api.db.nomics.world/v22/series/PROVIDER/DATASET/SERIES?observations=1}.
\end{center}
For each empirical result, record the endpoint, access date, series frequency,
provider, transformations, and treatment of missing values. Cite the original
provider and DBnomics as appropriate.

\section{Tested identifiers}

The following endpoints were executed successfully on 21 August 2026. A live
validation script is included in the repository. Identifiers can be revised
by providers; rerun it before using examples in a new project.

\begin{center}
\begin{tabular}{p{0.14\linewidth}p{0.23\linewidth}p{0.45\linewidth}}
\toprule
Frequency & Provider / dataset & Series code \\
\midrule
Monthly & ECB / EXR & \texttt{M.USD.EUR.SP00.A}: USD per euro spot average \\
Monthly & ECB / EXR & \texttt{M.GBP.EUR.SP00.A}: pound sterling per euro spot average \\
Quarterly & IMF / IFS & \texttt{Q.FR.NGDP\_R\_SA\_XDC}: French real GDP, seasonally adjusted \\
Annual & WB / WDI & \texttt{A-NY.GDP.MKTP.KD-FRA}: French GDP, constant 2015 US dollars \\
\bottomrule
\end{tabular}
\end{center}

The provider code for the World Bank example is \texttt{WB}. The intuitive
spelling \texttt{WorldBank} was checked and returns HTTP 404, so it must not
be used in code or teaching material.

\section{Python: download first, then transform}

The repository's Python examples require only the standard library. From a
clone, first validate every advertised identifier:

\begin{lstlisting}[language=bash]
python3 examples/python/validate_examples.py
\end{lstlisting}

To download one monthly series:
\begin{lstlisting}[language=bash]
python3 examples/python/download_series.py \
  ECB/EXR/M.USD.EUR.SP00.A --output /tmp/eurusd.csv
\end{lstlisting}

Several series can be written into the same long, tidy file:
\begin{lstlisting}[language=bash]
python3 examples/python/download_series.py \
  ECB/EXR/M.USD.EUR.SP00.A ECB/EXR/M.GBP.EUR.SP00.A \
  --output /tmp/exchange-rates.csv
\end{lstlisting}

The output has the columns \texttt{series}, \texttt{period}, and
\texttt{value}. Keeping a long format makes the identity of each series
explicit. The optional \texttt{--start} and \texttt{--end} arguments retain a
lexical range of standard annual, monthly, or quarterly period labels.

\texttt{annual\_growth.py} illustrates a transparent transformation: it
downloads annual French GDP and calculates a percent change only when both
adjacent observations are present. It leaves the first growth rate and values
next to missing observations blank. Do not silently interpolate, forward-fill,
or change frequency without a stated economic and statistical reason.

\section{MATLAB and a Dynare pre-processing step}

\texttt{examples/matlab/download\_series.m} uses \texttt{webread} and
\texttt{jsondecode} to create a table and, optionally, a CSV:

\begin{lstlisting}[language=Matlab]
T = download_series('IMF','IFS','Q.FR.NGDP_R_SA_XDC', ...
    'examples/output/france-real-gdp-quarterly.csv');
T.date = parse_dbnomics_period(T.period);
\end{lstlisting}

\texttt{parse\_dbnomics\_period.m} explicitly recognises only the annual
\texttt{YYYY}, monthly \texttt{YYYY-MM}, and quarterly \texttt{YYYY-Qq}
formats used in the examples. It fails for an unfamiliar convention, which is
safer than guessing dates. MATLAB represents unavailable numeric observations
as \texttt{NaN}; decide and document any subsequent treatment.

\texttt{download\_examples.m} produces local monthly, quarterly, and annual
CSV files and writes a two-column quarterly example suitable as a transparent
input stage before MATLAB/Dynare. A Dynare model should not assume an
unlabelled CSV: document the variable order, units, sample, transformations,
and any alignment or missing-data rule in the pre-processing script.

\section{Finding and checking a series}

\begin{enumerate}
\item Search \url{https://db.nomics.world/} by concept, country, and provider.
\item Read the title and dimensions; compare the frequency, unit, adjustment,
      and coverage with the intended model variable.
\item Copy the three identifiers from the page or API URL exactly.
\item Run a small request and inspect the first/last period and observation
      count. The repository validator does this for its examples.
\item Keep a human-readable data note alongside code: provider/dataset/series,
      URL, retrieval date, definitions, and all transformations.
\end{enumerate}

Do not publish large downloaded datasets when a deterministic downloader is
enough. If a project uses a restricted, changing, or manually curated input,
document the access constraint and publish only what redistribution permits.

\section{Checks before a release}

\begin{itemize}
\item Run the live Python validator and the downloader from a clean clone.
\item Confirm each API response returns the requested provider/dataset/series.
\item Check date parsing at the intended frequency and preserve missing values.
\item Rebuild this PDF from its \LaTeX{} source.
\item Keep only source, documentation, useful code, and the canonical PDF:
      no logs, caches, downloaded data, credentials, or local paths.
\end{itemize}

\end{document}
